%!TEX root = ../thesis.tex

\section{軟弱地面における床沈み込み現象}

軟弱地面とは，砂地や土壌，柔軟材料など，外力に応じて変形する床環境を指す．このような環境では，ロボットの足が接地した際に，床が弾性あるいは塑性的に変形し，足先が沈み込む現象が生じる．剛体床においては，足先の位置は接地後にほぼ不変であるのに対し，軟弱地面では接地後も床変形が進行する点が大きな違いである．

床沈み込み量は，床の剛性や減衰特性，ロボットの質量や接地速度など，複数の要因に依存する．そのため，沈み込み量は事前に正確に予測することが困難であり，歩行中に逐次的に発生する外乱としてロボットの運動に影響を与える．特に，片脚支持期において支持脚側の床が沈み込む場合，ロボット全体の姿勢や重心位置に顕著な影響を及ぼす．
